\documentclass[10pt,a4paper]{amsart}
\usepackage[margin=19mm]{geometry}
\usepackage{amsmath,amssymb,mathtools}
\usepackage[scr=rsfs,cal=euler]{mathalfa}
\usepackage{xfrac}
\usepackage[unicode,hidelinks,bookmarks=false]{hyperref}
\newcommand{\ie}{i.e.\ }
\newcommand{\cf}{cf.\ }
\newcommand{\resp}{respectively\ }
\newcommand{\Subsection}{\subsection}
\providecommand{\nobreakdash}{\nobreak\hbox{-}}
\setlength{\emergencystretch}{2em}
\multlinegap0pt
\usepackage{amssymb}
\usepackage{enumerate}
\usepackage{xcolor}
\newtheorem{proposition}{Proposition}
\newcommand{\Ab}{\mathbf A}
\newcommand{\C}{\mathbb C}
\newcommand{\Ef}{\mathcal E_\Phi}
\newcommand{\En}{\mathrm E(\Phi)}
\newcommand{\Fb}{\mathbf F}
\newcommand{\Gf}{\mathcal G_\Phi}
\newcommand{\Gn}{\mathrm G(\Phi)}
\newcommand{\N}{\mathbb N}
\newcommand{\R}{\mathbb R}
\newcommand{\ab}{\mathbf a}
\newcommand{\cE}{\mathcal E}
\newcommand{\cG}{\mathcal G}
\newcommand{\cH}{\mathcal H}
\newcommand{\cL}{\mathcal L}
\DeclareMathOperator{\curl}{curl}
\newcommand{\dd}{\mathop{}\!{\mathrm d}}
\newcommand{\ii}{\mathrm i}
\newcommand{\rG}{\mathrm G}
\title{Local strong magnetic fields and the Little-Parks effect}
\author{Ayman Kachmar, Mikael Sundqvist}
\date{Source: arXiv:2405.09099v2 (CC BY 4.0)}
\begin{document}
\maketitle

\noindent\emph{Source and licence.} Local strong magnetic fields and the Little-Parks effect. Version arXiv:2405.09099v2 (CC BY 4.0), distributed by arXiv under the Creative Commons Attribution 4.0 International licence, http://creativecommons.org/licenses/by/4.0/, as recorded in the arXiv licence metadata for this exact version and shown on the abstract page https://arxiv.org/abs/2405.09099v2. Full licence notice: \texttt{licenses/arxiv-2405.09099-CC-BY-4.0.txt}.\par\smallskip

\noindent\textit{Editorial scope.} Counted unit: the article's proposition prop:conv-min (source0.tex lines 519--527). The article's complete original proof of it is delivered with it (source0.tex lines 529--578, one \texttt{proof} environment of 50 lines). No mathematics is written, repaired or re-derived: the statement and the proof are the article's own lines, in the article's own notation, and no retained line is substituted or re-flowed.  Retained uncounted as the source-local context the proof needs: lines 219--221; lines 242--244; lines 246--248; lines 407--420; lines 454--457; lines 465--474; lines 471--473; lines 487--494.  Nothing was omitted from any retained span, so the package carries no omission record.  No retained line contains a citation, so the wrapper carries no bibliography and no bibliography entry.  No retained line contains a figure environment, an image-inclusion command or drawing code, so no image is delivered and the package declares no asset dependency.  Editorial formatting only, in addition: an AMS \texttt{amsart} preamble that restates the article's own declarations and macros used by the retained lines, namely the environments \texttt{proposition} on separate counters, together with the standard packages those lines need.  The retained environments print in this document as Proposition 1 in order of appearance, whereas the article prints the same items with the numbering of its own full document.  The complete native source of the article as distributed in the arXiv e-print for this exact version is bundled unchanged as \texttt{sources/159181/source0.tex}.
\begin{equation}\label{eq:def-energy}
\En=\inf\{\Ef(\psi,\Ab)\colon (\psi,\Ab)\in\cH\}.
\end{equation}

\begin{equation}\label{eq:def-energy0}
\Gn=\inf\{\Gf(u,\Ab)\colon (u,\Ab)\in \cH_0\}.
\end{equation}

\begin{equation}\label{eq:def-U-1}
U_n(x)=\exp\left(\ii n\int_{\ell(x_*,x)}\Fb\cdot\dd \mathrm r\right),
\end{equation}

\begin{proposition}\label{prop:GL-eff}
For all $\Phi>0$, we have
\[
    \rG(\Phi+1)=\Gn,
\]
where  $\rG(\cdot)$ is introduced in \eqref{eq:def-energy0}. Moreover, with $n=\lfloor \Phi\rfloor$ and $\Phi_0=\{\Phi\}$, we have
\begin{multline*}
    (u_*,\Ab_*)\text{ is a minimizer of }\cG_{\Phi_0}
    \iff \\
    (u,\Ab)
    =
    \cL_{\Phi_0,n}(u_*,\Ab_*)\text{ is a minimizer of }\Gf.
\end{multline*}
\end{proposition}

\begin{proposition}\label{prop:ub}
Let $\En$ and $\Gn$ be the ground state energies introduced in \eqref{eq:def-energy} and \eqref{eq:def-energy0} respectively. Then, for all $\Phi\geq 0$, we have
\[ \En\leq \Gn\leq 0.\]
\end{proposition}

\begin{proposition}\label{prop:est-min}
There exists a positive constant $C$ such that, if $\Phi>0$ and $(\psi_\Phi,\Ab_\Phi)\in\cH$ is a minimizer  of $\Ef$,  then
\[\begin{gathered}
\|\psi\|_\infty\leq 1,\quad \|(-\ii\nabla-\Phi\Ab)\psi\|_2\leq \kappa\|\psi\|_2,\quad \|\curl(\Ab-\Fb)\|_2\leq \frac{\kappa}{\Phi}\|\psi\|_2,
\end{gathered}\]
and
\begin{equation}\label{eq:curl-div}
\|\Ab-\Fb\|_{H^2(\Omega)}\leq \frac{C\kappa}{\Phi}\|\psi\|_2.
\end{equation}
\end{proposition}

\begin{equation}
\|\Ab-\Fb\|_{H^2(\Omega)}\leq \frac{C\kappa}{\Phi}\|\psi\|_2.
\end{equation}

\begin{proposition}\label{prop:conv-density}
Let  $(\psi_\Phi,\Ab_\Phi)_{\Phi>0}\subset\cH$ be a family of configurations such that $(\psi_\Phi,\Ab_\Phi)$ is a minimizer of $\Ef$ for every $\Phi>0$. Then, there exists a configuration $(\rho_*,\ab_*)\in\cH_0$ and a sequence $(\Phi_n)$ such that $\Phi_n\to+\infty$ and
\[
\begin{gathered}
|\psi_{\Phi_n}|\to \rho_*\text{ weakly in }H^1(\Omega;\R)\text{ and strongly in }L^2(\Omega;\R),\\
\Phi_n(\Ab_{\Phi_n}-\Fb)\to \ab_*\text{ weakly in }H^2(\Omega;\R^2) \text{ and strongly in }H^1(\Omega;\R^2).
\end{gathered}\]
\end{proposition}

\begin{proposition}\label{prop:conv-min}
Let  $(\psi_\Phi,\Ab_\Phi)_{\Phi>0}\subset\cH$ be a family of configurations such that $(\psi_\Phi,\Ab_\Phi)$ is a minimizer of $\Ef$ for every $\Phi>0$. Then, there exists a sequence $(\Phi_n)$ and a function $u_*\in H^1(\Omega;\C)$ such that $\Phi_n\to+\infty$ and
\[\begin{gathered}
\cE_{\Phi_n}(\psi_{\Phi_n},\Ab_{\Phi_n})\geq \rG(\Phi_n)+o(1),\\
 U_{\lfloor\Phi_n\rfloor}^{-1}\psi_{\Phi_n}\to u_*\text{ in }H^1(\Omega_0;\C),\\
 u_*=0\text{ on }\omega,
 \end{gathered}\]
 where, for $k\in\N$, $U_k$ is the function introduced in \eqref{eq:def-U-1}.
\end{proposition}

\begin{proof}
The proof relies on a compactness argument and the extraction of a subsequence. To simplify the exposition, we will skip the reference to the subsequence.

For $\Phi>0$, we write $\Phi=\Phi_0+n$ where $n=\lfloor \Phi\rfloor$ and $\Phi_0=\{\Phi\}\in [0,1)$. Without loss of generality, we will assume that $\Phi_0$ is a fixed constant. Starting from $(\psi,\Ab)=(\psi_\Phi,\Ab_\Phi)$, we put
\[ \rho=|\psi|,\quad u=u_\Phi:=U_n^{-1}\psi,\quad \ab=\ab_\Phi:=\Phi_0\Fb+\Phi(\Ab-\Fb), \]
and we observe by \eqref{eq:curl-div} that
\[ \|\ab\|_{H^2(\Omega)}\leq C_0,\]
where $C_0$ is a positive constant, independent of $\Phi$.  Moreover, note that $u$ is defined on $\Omega_0$,  $|u|=\rho|_{\Omega_0}$, and the  following identities hold in $\Omega_0$,
\[\begin{gathered}
(-\ii\nabla-\ab)u=U_n^{-1}(-\ii\nabla-\Phi\Ab)\psi,\\
(-\ii\nabla-\ab)^2u=U_n^{-1}(-\ii\nabla-\Phi\Ab)^2\psi=\kappa^2(1-|u|^2)u.
\end{gathered}\]
Consequently, with
\begin{equation}\label{eq:def-Om-ep}
0<\varepsilon<\varepsilon_0:=\mathrm{dist}(\partial\omega,\partial\Omega),\quad \Omega_\varepsilon:=\{x\in\Omega_0\colon \mathrm{dist}(x,\partial\omega)>\varepsilon\},
\end{equation} we have  by Proposition~\ref{prop:est-min} and elliptic estimates,
\[\|u\|_{H^2(\Omega_\varepsilon)}\leq C_\varepsilon, \]
where $C_\varepsilon$ depends on $\varepsilon$ but is independent of $\Phi$.


Cantor's diagonal argument yields a sequence  and a  function $u_*\in H^2_{\rm loc}(\Omega_0;\C)$ such that,  for $0<\varepsilon<\varepsilon_0$,
\begin{equation}\label{eq:proof-conv}
 u\to u_*\text{ weakly in }H^2(\Omega_\varepsilon;\C) \mbox{ and strongly in }H^1(\Omega_\varepsilon;\C),\end{equation}
and by Propositions~\ref{prop:est-min} and ~\ref{prop:conv-density},
\begin{equation}\label{eq:proof-conv*}
\begin{gathered}
\|u_*\|_\infty\leq 1,\\
\rho\to \rho_*\text{ in }L^2(\Omega;\R)\text{ with }\rho_*|_{\omega}=0,\\
\ab\to \Phi_0\Fb+\ab_*\text{ in }H^1(\Omega;\R^2).
\end{gathered}
\end{equation}
Consequently,  $|u_*|=\rho_*|_{\Omega_0}$.  Moreover, by monotone convergence
\[\int_{\Omega_0}|(-\ii\nabla-\Phi_0\Ab_*)u_*|^2\dd x=\lim_{\varepsilon\to0_+}\int_{\Omega_\varepsilon}|(-\ii\nabla-\Phi_0\Ab_*)u_*|^2\dd x, \]
where
\[\Ab_*=\begin{cases}\Fb+\Phi_0^{-1}\ab_*&\text{if }\Phi_0>0,\\
\Fb&\text{if }\Phi_0=0,\end{cases}\]
and we have by \eqref{eq:proof-conv}-\eqref{eq:proof-conv*} and Proposition~\ref{prop:est-min},
\begin{multline*}\forall\varepsilon\in(0,\varepsilon_0),\quad \int_{\Omega_\varepsilon}|(-\ii\nabla-\Phi_0\Ab_*)u_*|^2\dd x\\
\leq \limsup_{\Phi\to+\infty}\int_{\Omega}|(-\ii\nabla-\Phi\Ab)\psi|^2\dd x\leq \kappa^2|\Omega|,\end{multline*}
hence it follows that
\[\int_{\Omega_0}|(-\ii\nabla-\Phi_0\Ab_*)u_*|^2\dd x \leq \kappa^2|\Omega|.\]
 Thus, we have $(u_*,\Ab_*)\in\cH_0$ and, for $0<\varepsilon<\varepsilon_0$,
\[\begin{aligned}
\En&=\Ef(\psi,\Ab)\\
&\geq \int_{\Omega_\varepsilon}\Bigl(|(-\ii\nabla-\ab)u|^2+\frac{\kappa^2}2|u|^4\Bigr)\dd x-\kappa^2\int_{\Omega}\rho^2\dd x+\int_\Omega|\curl\ab|^2\dd x .
\end{aligned}\]
Sending $\Phi\to+\infty$ (along the sequence $(\Phi_n)$) and then $\varepsilon\to0$, we get by monotone convergence
\[\liminf_{\Phi_n\to+\infty}\En\geq \cG_{\Phi_0}(u_*,\Ab_*)\geq \rG(\Phi_0).\]
To finish the proof, we use Propositions~\ref{prop:ub} and \ref{prop:GL-eff}, and this yields that $(u_*,\Ab_*)$ is a minimizer of $\cG_{\Phi_0}$.
\end{proof}

\end{document}
